\documentclass[11pt]{article}

\usepackage[legalpaper,margin=0.5in]{geometry}
\usepackage{titlesec}
\usepackage{enumitem}
\usepackage{hyperref}

\pagenumbering{gobble}

\titleformat{\section}
{\large\bfseries}
{}{0em}{}

\begin{document}

\begin{center}
{\LARGE \textbf{Derrickson Ymbon}}\\[4pt]
Department of Mathematics\\
North-Eastern Hill University\\
Shillong, India\\[4pt]
Email: \href{mailto:youremail@email.com}{derrixonymbon202@email.com}
\end{center}

\vspace{0.4cm}

\section*{Research Interests}

Quiver Algebra, Representation Theory and Category Theory.

\section*{Education}


\noindent \textbf{M.Sc. in Mathematics} \hfill 2024 - \\
CGPA: 6.5\\
North-Eastern Hill University

\vspace{0.2cm}

\noindent \textbf{B.Sc. in Mathematics} \hfill 2021 - 2024\\
CGPA: 5.5\\
North-Eastern Hill University 



\section*{Relevant Courses}
Graduate Algebra, Analysis, Linear Algebra, Topology, Differential Geometry. 

% \begin{itemize}
%     \item \textbf{Graduate Algebra}(Approx. 60 Hours) -  
%     \item \textbf{Analysis}(Approx. 60 Hours) -
%     \item \textbf{Linear Algebra}(Approx. 60 Hours) -
% \end{itemize}

\section*{Projects}

\begin{itemize}[leftmargin=*, itemsep=0.6em]

\item \textbf{Master's Project: Surface Triangulations and the Structures of Gentle and Brauer Graph Algebras} \\
\textit{(February 2026 -- present)} \\
Supervisors: Dr.~Ardeline Mary Buhphang and Prof.~Thomas Br\"{u}stle, University of Sherbrooke, Canada.

\vspace{3mm}

Study of the relationship between triangulated surfaces and the associated classes of gentle algebras and Brauer graph algebras, focusing on their combinatorial and representation-theoretic structures.

\item \textbf{Reading Project: Homological Algebra in Representation Theory} \\
\textit{(January 2026 -- present)} \hfill Mentor: Dr.~Ardeline Mary Buhphang \\

Ongoing guided reading, the project follows Henning Krause's
\textit{Homological Theory of Representations} together with C.~Weibel's
\textit{An Introduction to Homological Algebra}. 

Topics studied include projective and injective modules, projective and injective
resolutions, derived functors such as $\mathrm{Ext}$ and $\mathrm{Tor}$, chain
complexes and chain homotopies, and the construction of homotopy and derived
categories. Particular emphasis is placed on the role of these homological tools
in the study of module categories of finite-dimensional algebras and in the
representation theory of quivers.

\item \textbf{Reading Project: Quiver Representations and Auslander--Reiten Theory} \\
\textit{(February 2026 -- present)} \\ \hfill Mentors: Prof.~Thomas Br\"ustle (University of Sherbrooke), Dr.~Ardeline Mary Buhphang (NEHU) \\

Study of quiver representations based on Ralf Schiffler's \textit{Quiver Representations}, covering Gabriel's theorem, indecomposable representations, and an introduction to Auslander--Reiten theory.

\end{itemize}

\section*{Computer Skills}
\begin{itemize}
    \item Progamming Language: $C$, Python.
    \item Typesetting: \LaTeX
\end{itemize}









\end{document}